\documentclass[11pt,a4paper]{article}
\usepackage[italian]{babel}
\usepackage[margin=2.5cm]{geometry}
\usepackage{parskip}
\usepackage{amsmath, amssymb, amsthm}
\usepackage{booktabs}
\usepackage{array}
\usepackage{xcolor}
\usepackage{needspace}
\usepackage{enumitem}
\usepackage{listings}
\usepackage{tikz}
\usepackage[most]{tcolorbox}
\usepackage[hidelinks]{hyperref}
\usetikzlibrary{decorations.pathreplacing, shapes.geometric, arrows.meta, positioning}

% Niente vedove/orfane: i paragrafi non lasciano righe isolate a fine/inizio pagina
\widowpenalty=10000
\clubpenalty=10000
\linespread{0.97}

% Elenchi più compatti
\setlist{itemsep=2pt, parsep=0pt, topsep=4pt, partopsep=0pt}

% --- Riquadri con bordo nero, senza numero (si spezzano tra due pagine) ---
% Uso: \begin{definizione} ... \end{definizione}
%      \begin{teorema}[Nome] ... \end{teorema}  (il nome è facoltativo)
\tcbset{nota/.style={enhanced jigsaw, breakable, colback=white,
  colframe=black, boxrule=0.5pt, sharp corners}}
\NewTColorBox{definizione}{}{nota,
  before upper={\textbf{Definizione.}\ }}
\NewTColorBox{teorema}{o}{nota,
  before upper={\textbf{Teorema\IfValueT{#1}{ (#1)}.}\ }}
\NewTColorBox{proposizione}{o}{nota,
  before upper={\textbf{Proposizione\IfValueT{#1}{ (#1)}.}\ }}
\NewTColorBox{proprieta}{o}{nota,
  before upper={\textbf{Proprietà\IfValueT{#1}{ (#1)}.}\ }}
\NewTColorBox{assioma}{o}{nota, boxrule=1pt,
  before upper={\textbf{Assioma\IfValueT{#1}{ (#1)}.}\ }}

% --- Ambienti semplici (senza riquadro) ---
% Il titolo sta in cima al blocco, anche se il contenuto inizia con un riquadro o
% con codice; e non resta mai da solo in fondo a una pagina.
% Uso: \begin{esempio}[Titolo facoltativo] ... \end{esempio}
\newcommand{\newrunin}[2]{%
  \NewDocumentEnvironment{#1}{o}{%
    \par\Needspace{4\baselineskip}\addvspace{\medskipamount}\noindent
    \textbf{#2}\IfValueT{##1}{ (##1)}\textbf{.}\ \ignorespaces}%
  {\par\addvspace{\medskipamount}}}
\newrunin{esempio}{Esempio}
\newrunin{esercizio}{Esercizio}
\NewTColorBox{osservazione}{}{nota,
  before upper={\textbf{Osservazione.}\ }}

% V e F nelle tavole di verità
\newcommand{\V}{\textsf{V}}
\newcommand{\F}{\textsf{F}}

% --- Codice C ---
% Uso: \begin{codice} ... \end{codice}      codice C numerato
%      \begin{comando} ... \end{comando}    comandi da terminale
%      \begin{risultato} ... \end{risultato} output di un programma
%      \lstinline|printf("ciao");|           codice dentro al testo
\lstdefinestyle{c}{
  language=C,
  basicstyle=\ttfamily\small,
  keywordstyle=\color{blue}\bfseries,
  commentstyle=\color{gray}\itshape,
  stringstyle=\color{red!70!black},
  numbers=left,
  numberstyle=\tiny\color{gray},
  numbersep=8pt,
  columns=flexible,
  keepspaces=true,
  breaklines=true,
  showstringspaces=false,
  tabsize=4,
  morekeywords={include, define},
  % lettere accentate nei commenti
  literate={à}{{\`a}}1 {è}{{\`e}}1 {é}{{\'e}}1 {ì}{{\`i}}1
           {ò}{{\`o}}1 {ù}{{\`u}}1 {È}{{\`E}}1 {À}{{\`A}}1
}
\lstset{style=c}
\newtcblisting{codice}{listing only, nota, left=6mm,
  listing options={style=c}}
\newtcblisting{comando}{listing only, nota,
  listing options={basicstyle=\ttfamily\small, numbers=none, language={},
    columns=flexible, keepspaces=true, breaklines=true}}
\newtcblisting{risultato}{listing only, enhanced jigsaw, breakable,
  colback=black!5, colframe=black, boxrule=0.5pt, sharp corners,
  listing options={basicstyle=\ttfamily\small, numbers=none, language={},
    columns=flexible, keepspaces=true, breaklines=true}}

% --- Diagrammi di flusso (simboli standard) ---
\tikzset{
  terminale/.style = {ellipse, draw, minimum width=2.5cm, minimum height=0.9cm},
  io/.style        = {trapezium, trapezium left angle=70, trapezium right angle=110,
                      draw, minimum width=2.5cm, minimum height=0.9cm},
  processo/.style  = {rectangle, draw, minimum width=2.5cm, minimum height=0.9cm},
  decisione/.style = {diamond, draw, aspect=2, minimum width=2.5cm},
  freccia/.style   = {thick, -{Stealth}}
}

\title{Reti logiche\\[0.4em]\large Appunti}
\author{Marco Cerbella}
\date{Università degli Studi di Perugia -- A.A. 2026/2027\\ Aggiornato al \today}

\begin{document}
\maketitle
\setcounter{tocdepth}{1}
\tableofcontents
\clearpage

\section[Algebra di Boole e algebra di commutazione (Lezione 1 -- 21 settembre 2026)]{Algebra di Boole e algebra di commutazione}
\noindent\emph{Lezione 1 -- 21 settembre 2026}
\subsection*{Cenno storico}

L'algebra di Boole fu introdotta da George Boole (1815--1864) nell'Ottocento per
dimostrare affermazioni logiche svolgendo calcoli algebrici: per Boole la
logica va associata alla matematica e non alla filosofia. Nel 1936 Claude
Shannon (1916--2001), studiando al MIT un circuito di oltre cento relè
collegato all'analizzatore differenziale di Bush, intuì che l'algebra di Boole
è un modello formale per l'\emph{analisi} e la \emph{sintesi} di tali sistemi:
due interruttori in \emph{serie} (\emph{parallelo}) si descrivono con AND (OR),
mentre il NOT si realizza con il contatto posteriore di un relè. Dalla sua tesi
(1938) nasce l'algebra di commutazione, base delle reti logiche.

\subsection{Algebra di Boole}

\begin{definizione}
Un'\emph{algebra di Boole} è una struttura algebrica
\[
  \mathcal{B} = \langle A, +, \cdot, ', 0, 1 \rangle
\]
dove:
\begin{itemize}
    \item $A$ è l'insieme degli elementi su cui l'algebra opera (\emph{supporto});
    \item $+$ è l'operatore di disgiunzione (OR);
    \item $\cdot$ è l'operatore di congiunzione (AND);
    \item $'$ è l'operatore unario di complementazione (NOT);
    \item $0$ è l'elemento neutro rispetto a $+$, $1$ è l'elemento neutro
          rispetto a $\cdot$;
\end{itemize}
e valgono i seguenti assiomi (\emph{postulati di Huntington}).
\end{definizione}

\subsubsection{Postulati di Huntington}

\begin{enumerate}
    \item \textbf{Commutatività}: $\forall a, b \in A$:
          $a + b = b + a$ \quad e \quad $a \cdot b = b \cdot a$.
    \item \textbf{Elementi neutri}: esistono $0$ e $1$ tali che, per ogni
          $a \in A$, $a + 0 = a$ e $a \cdot 1 = a$.
    \item \textbf{Distributività} (di ciascuna operazione rispetto all'altra):
          $\forall a, b, c \in A$
          \begin{align*}
            a + (b \cdot c) &= (a + b) \cdot (a + c)\\
            a \cdot (b + c) &= (a \cdot b) + (a \cdot c)
          \end{align*}
    \item \textbf{Complemento}: per ogni $a \in A$ esiste ed è unico $a'$ tale
          che $a + a' = 1$ e $a \cdot a' = 0$.
\end{enumerate}

\begin{teorema}[Principio di dualità]
Ogni proposizione derivata dagli assiomi di Huntington resta valida scambiando
ovunque tra loro le operazioni $+$ e $\cdot$ e gli elementi neutri $0$ e $1$.
\end{teorema}

\begin{osservazione}
In particolare se $\langle A, +, \cdot, ', 0, 1 \rangle$ è un'algebra di Boole,
lo è anche $\langle A, \cdot, +, ', 1, 0 \rangle$.
\end{osservazione}

\subsubsection{Esempi}

\begin{esempio}[Algebra delle classi]
Sia $A$ un insieme e $\mathcal{P}(A)$ l'insieme delle sue parti (tutti i
sottoinsiemi di $A$). La struttura
$\langle \mathcal{P}(A), \cup, \cap, \overline{\phantom{x}}, \emptyset, A \rangle$
è un'algebra di Boole:
\begin{itemize}
    \item $\cup$ e $\cap$ sono commutative;
    \item $\emptyset$ è neutro rispetto a $\cup$ e $A$ rispetto a $\cap$:
          per ogni $X \in \mathcal{P}(A)$, $X \cup \emptyset = X$ e
          $X \cap A = X$;
    \item distributività e unicità del complemento si verificano con i
          diagrammi di Venn.
\end{itemize}
È un'algebra di Boole \emph{chiusa}: $\forall X, Y \in \mathcal{P}(A)$,
$X \cap Y$, $X \cup Y$ e $\overline{X} = A \setminus X$ sono sottoinsiemi di $A$,
quindi elementi di $\mathcal{P}(A)$.
\end{esempio}

\begin{teorema}[Rappresentazione di Stone]
Per ogni algebra di Boole finita $\mathcal{A}$ esiste un insieme finito $X$ tale
che $\mathcal{A}$ è isomorfa a
$\langle \mathcal{P}(X), \cup, \cap, \overline{\phantom{x}}, \emptyset, X \rangle$.
\end{teorema}

Conseguenze:
\begin{itemize}
    \item la cardinalità del supporto di un'algebra di Boole finita è sempre una
          potenza di $2$;
    \item per dimostrare che le algebre di Boole finite godono di una proprietà
          basta dimostrarla per l'algebra delle classi (ad esempio con i
          diagrammi di Venn).
\end{itemize}

\begin{esempio}[Altre algebre di Boole]
$\langle \{0,1\}, \min, \max, 1-x, 0, 1 \rangle$ è un'algebra di Boole; così
come quella ottenuta scambiando le operazioni come nell'osservazione sopra.
\end{esempio}

\subsubsection{Variabili booleane}

Sia $\mathcal{A} = \langle A, +, \cdot, ', 0, 1 \rangle$ un'algebra di Boole. Una
\emph{variabile booleana} è un simbolo che indica un qualsiasi elemento del
supporto $A$; il suo valore è un elemento di $A$. Le variabili si indicano con
le ultime lettere dell'alfabeto ($x, y, \dots$), gli elementi del supporto con
le prime ($a, b, \dots$) e con i simboli $0$ e $1$.

\subsubsection{Proprietà dell'algebra di Boole}

\begin{proprieta}
Per ogni $x, y, z$ valgono:
\begin{itemize}
    \item \textbf{associatività}: $x + (y + z) = (x + y) + z$ \ e \
          $x \cdot (y \cdot z) = (x \cdot y) \cdot z$;
    \item \textbf{idempotenza}: $x + x = x$ \ e \ $x \cdot x = x$;
    \item \textbf{elemento nullo}: $x + 1 = 1$ \ e \ $x \cdot 0 = 0$;
    \item \textbf{unicità del complemento}: $x'$ è unico;
    \item \textbf{assorbimento}: $x + x \cdot y = x$ \ e \
          $x \cdot (x + y) = x$;
    \item \textbf{semplificazione}: $x + x' \cdot y = x + y$ \ e \
          $x \cdot (x' + y) = x \cdot y$;
    \item \textbf{involuzione}: $(x')' = x$;
    \item \textbf{leggi di De Morgan}: $(x + y)' = x' \cdot y'$ \ e \
          $(x \cdot y)' = x' + y'$;
    \item \textbf{consenso}:
          $x \cdot y + x' \cdot z + y \cdot z = x \cdot y + x' \cdot z$ \ e \
          $(x + y) \cdot (x' + z) \cdot (y + z) = (x + y) \cdot (x' + z)$.
\end{itemize}
\end{proprieta}

Ciascuna si può verificare con i diagrammi di Venn, oppure dimostrare dagli
assiomi.

\begin{esempio}[Dimostrazione dell'idempotenza]
\begin{align*}
x \cdot x &= (x \cdot x) + 0            && \text{elemento neutro}\\
          &= (x \cdot x) + (x \cdot x') && \text{complemento}\\
          &= x \cdot (x + x')           && \text{distributiva}\\
          &= x \cdot 1                  && \text{complemento}\\
          &= x                          && \text{elemento neutro}
\end{align*}
L'idempotenza di $+$ ($x + x = x$) segue per dualità.
\end{esempio}

\subsection{Algebra di commutazione}

\begin{definizione}[Algebra di commutazione (\emph{switching algebra})]
È l'algebra di Boole il cui supporto è costituito dai soli valori $\{0, 1\}$:
\[
  \mathcal{B} = \langle \{0,1\}, +, \cdot, ', 0, 1 \rangle .
\]
Ogni variabile può assumere solo i valori $0$ e $1$.
\end{definizione}

È l'algebra usata da Shannon per l'analisi dei circuiti. I tre operatori sono
definiti dalle seguenti tabelle di verità.

\begin{center}
\begin{tabular}{c|cc}
$+$ & 0 & 1 \\ \midrule
0 & 0 & 1 \\
1 & 1 & 1
\end{tabular}
\qquad\qquad
\begin{tabular}{c|cc}
$\cdot$ & 0 & 1 \\ \midrule
0 & 0 & 0 \\
1 & 0 & 1
\end{tabular}
\qquad\qquad
\begin{tabular}{c|c}
$x$ & $x'$ \\ \midrule
0 & 1 \\
1 & 0
\end{tabular}
\end{center}

\subsubsection{OR esclusivo}

L'\emph{OR esclusivo} (XOR) è l'operatore $\oplus$ definito da
\[
  a \oplus b = a' b + a b' ,
\]
cioè vale $1$ quando $a$ e $b$ sono diversi.

\begin{center}
\begin{tabular}{c|cc}
$\oplus$ & 0 & 1 \\ \midrule
0 & 0 & 1 \\
1 & 1 & 0
\end{tabular}
\end{center}

Dalla definizione, applicando le proprietà dell'algebra di Boole, si ricava:
\begin{align}
x \oplus 1 &= x' & x \oplus 0 &= x \\
x \oplus x &= 0 & x \oplus x' &= 1 \\
x' \oplus y &= x \oplus y' = (x \oplus y)'
\end{align}

\subsection*{Esercizi per casa}

\begin{esercizio}
Usare i diagrammi di Venn per dimostrare le proprietà (associatività,
idempotenza, \dots) dell'algebra booleana su supporto finito.
\end{esercizio}
\begin{esercizio}
Dimostrare che l'algebra di commutazione è un'algebra di Boole, cioè verificare
i postulati di Huntington.
\end{esercizio}
\begin{esercizio}
Dimostrare, per l'algebra di commutazione, associatività, idempotenza, elemento
nullo, unicità del complemento, assorbimento, semplificazione, involuzione, De
Morgan e consenso.
\end{esercizio}

\section[Funzioni ed espressioni booleane (Lezione 2 -- 6 ottobre 2026)]{Funzioni ed espressioni booleane}
\noindent\emph{Lezione 2 -- 6 ottobre 2026}
\subsection{Funzioni booleane}

\begin{definizione}[Funzione booleana]
Sia $\mathcal{B} = \langle A = \{0,1\}, +, \cdot, ', 0, 1 \rangle$ un'algebra di
commutazione. Una \emph{funzione booleana} a $n$ variabili è una relazione
\[
  f \colon A^n \to A
\]
che mette in corrispondenza gli elementi del dominio $A^n$ con gli elementi del
codominio $A$.
\end{definizione}

Dal punto di vista operativo una funzione booleana si definisce con le quattro
regole seguenti.
\begin{enumerate}
    \item $\forall a \in A$, la \emph{funzione costante}
          $f(x_1, \dots, x_n) = a$ è una funzione booleana a $n$ variabili.
    \item $\forall x_i \in \{x_1, \dots, x_n\}$, la \emph{funzione proiezione}
          $f(x_1, \dots, x_n) = x_i$ è una funzione booleana a $n$ variabili.
    \item Se $g$ e $h$ sono funzioni booleane a $n$ variabili, lo sono anche
          $g + h$, $g \cdot h$ e $g'$, definite da
          \begin{align*}
            (g + h)(x_1, \dots, x_n) &= g(x_1, \dots, x_n) + h(x_1, \dots, x_n)\\
            (g \cdot h)(x_1, \dots, x_n) &= g(x_1, \dots, x_n) \cdot h(x_1, \dots, x_n)\\
            (g')(x_1, \dots, x_n) &= \bigl(g(x_1, \dots, x_n)\bigr)'
          \end{align*}
    \item Non esistono altre funzioni booleane oltre a quelle ottenute
          dall'applicazione iterativa delle regole precedenti.
\end{enumerate}

\subsubsection{Tabelle di verità}

In un'algebra di commutazione $A = \{0,1\}$, quindi le funzioni booleane hanno
la forma
\[
  f \colon \{0,1\}^n \to \{0,1\}.
\]
Il dominio di $f$ ha $2^n$ elementi e il codominio ne ha $2$. La funzione può
essere descritta dalla sua \emph{tabella di verità} $T_f$, che contiene una riga
per ciascuna configurazione delle $n$ variabili in ingresso, con il relativo
valore di uscita di $f$.

\begin{esempio}
Funzione booleana sulle variabili $x, y, z$ che vale $1$ se e solo se
\emph{esattamente due} variabili valgono $1$.
\begin{center}
\begin{tabular}{ccc|c}
$x$ & $y$ & $z$ & $f(x,y,z)$ \\ \midrule
0 & 0 & 0 & 0 \\
0 & 0 & 1 & 0 \\
0 & 1 & 0 & 0 \\
0 & 1 & 1 & 1 \\
1 & 0 & 0 & 0 \\
1 & 0 & 1 & 1 \\
1 & 1 & 0 & 1 \\
1 & 1 & 1 & 0
\end{tabular}
\end{center}
\end{esempio}

\begin{proposizione}
Il numero di funzioni booleane (su un'algebra di commutazione) di $n$ variabili
è $2^{2^n}$.
\end{proposizione}

Infatti, date $n$ variabili si hanno $2^n$ configurazioni (tante quanti gli
elementi di $\{0,1\}^n$) e ogni funzione è definita dall'insieme delle
configurazioni per cui assume valore $1$. Il numero di funzioni è quindi la
cardinalità dell'insieme delle parti di $\{0,1\}^n$, cioè $2^{2^n}$.

\subsubsection{Esempi}

\begin{esempio}[Funzioni di una variabile]
Alcune di queste hanno particolare importanza dal punto di vista della
realizzazione.
\begin{center}
\begin{tabular}{c|cccc}
$x$ & $f(x)$ & $f(x)$ & $f(x)$ & $f(x)$ \\ \midrule
0 & 0 & 0 & 1 & 1 \\
1 & 0 & 1 & 0 & 1 \\ \midrule
  & \footnotesize costante $0$ & \footnotesize proiezione di $x$
  & \footnotesize negazione di $x$ & \footnotesize costante $1$
\end{tabular}
\end{center}
\end{esempio}

\begin{esempio}[Funzioni di due variabili (I)]
\begin{center}
\small
\begin{tabular}{cc|cccccc}
$x$ & $y$ & \multicolumn{6}{c}{$f(x,y)$} \\ \midrule
0 & 0 & 0 & 1 & 0 & 0 & 1 & 1 \\
0 & 1 & 0 & 1 & 0 & 1 & 1 & 0 \\
1 & 0 & 0 & 1 & 1 & 0 & 0 & 1 \\
1 & 1 & 0 & 1 & 1 & 1 & 0 & 0 \\ \midrule
\multicolumn{2}{c|}{}
  & \footnotesize costante $0$ & \footnotesize costante $1$
  & \footnotesize proiez.\ $x$ & \footnotesize proiez.\ $y$
  & \footnotesize negaz.\ $x$ & \footnotesize negaz.\ $y$
\end{tabular}
\end{center}
\end{esempio}

\begin{esempio}[Funzioni di due variabili (II)]
\begin{center}
\small
\begin{tabular}{cc|cccccc}
$x$ & $y$ & \multicolumn{6}{c}{$f(x,y)$} \\ \midrule
0 & 0 & 0 & 0 & 1 & 1 & 0 & 1 \\
0 & 1 & 0 & 1 & 1 & 0 & 1 & 0 \\
1 & 0 & 0 & 1 & 1 & 0 & 1 & 0 \\
1 & 1 & 1 & 1 & 0 & 0 & 0 & 1 \\ \midrule
\multicolumn{2}{c|}{}
  & \footnotesize AND & \footnotesize OR & \footnotesize NAND
  & \footnotesize NOR & \footnotesize XOR & \footnotesize EQUIV.
\end{tabular}
\end{center}
\end{esempio}

\subsubsection{Funzioni non completamente specificate}

Talvolta una funzione $f$ viene specificata solo in parte e si dice
\emph{non completamente specificata}. Ci sono due motivi principali:
\begin{enumerate}
    \item la natura del problema garantisce che una certa combinazione delle
          variabili di ingresso non possa presentarsi;
    \item non ha importanza quale valore assuma la funzione per una certa
          combinazione delle variabili di ingresso.
\end{enumerate}
Nel primo caso si dice che la funzione \emph{non è controllabile} in
corrispondenza di quella combinazione di ingressi, nel secondo che
\emph{non è osservabile}.

\begin{esempio}
Sia $f$ una funzione di due variabili $x_1, x_2$ tale che
\[
  \bigl(f(x_1,x_2) = 1 \iff x_1 = x_2 = 0\bigr) \ \wedge\
  \bigl(f(x_1,x_2) = 0 \iff x_1 \neq x_2\bigr).
\]
Il valore di $f(1,1)$ non è specificato (nella tabella si indica con $-$).
\begin{center}
\begin{tabular}{cc|c}
$x_1$ & $x_2$ & $f(x_1,x_2)$ \\ \midrule
0 & 0 & 1 \\
0 & 1 & 0 \\
1 & 0 & 0 \\
1 & 1 & $-$
\end{tabular}
\end{center}
\end{esempio}

\subsection{Espressioni booleane}

\begin{definizione}[Espressione booleana]
Un'\emph{espressione booleana} definita su un insieme di variabili booleane
$X = \{x_1, x_2, \dots, x_n\}$ è definita dalle regole seguenti:
\begin{itemize}
    \item gli elementi del supporto sono espressioni;
    \item le variabili $x_1, x_2, \dots, x_n$ sono espressioni;
    \item se $h, g$ sono espressioni booleane, lo sono anche $h + g$,
          $h \cdot g$ e $h'$;
    \item non esistono altre espressioni oltre a quelle ottenute applicando
          iterativamente le regole precedenti.
\end{itemize}
\end{definizione}

\subsubsection{Notazioni e convenzioni}

Fissato l'insieme dei simboli usati per le variabili, si può eliminare il
simbolo di prodotto logico $\cdot$ nei termini prodotto. Per le proprietà
dell'algebra booleana si possono inoltre omettere alcune parentesi, ricordando
che il prodotto logico ha \emph{priorità} sulla somma logica.

\begin{esempio}
\begin{itemize}
    \item $x + (y \cdot z)$ si scrive $x + yz$ (si elimina il simbolo $\cdot$
          nel termine prodotto $y \cdot z$ e si usa la priorità di $\cdot$
          rispetto a $+$);
    \item per la proprietà associativa, $x + (y + z \cdot x)$ si scrive
          $x + y + zx$ (si eliminano le parentesi per l'associatività e il
          simbolo $\cdot$ nel termine prodotto $z \cdot x$).
\end{itemize}
\end{esempio}

\subsubsection{Letterali}

Un \emph{letterale} è un'occorrenza di una variabile (in forma naturale o
negata) all'interno di un'espressione.

\begin{esempio}
L'espressione $(x' + y') \cdot (x + x')$ contiene $4$ letterali:
\begin{itemize}
    \item un'occorrenza della variabile $x$ in forma naturale;
    \item due occorrenze della variabile $x$ in forma negata ($x'$);
    \item un'occorrenza della variabile $y$ in forma negata ($y'$).
\end{itemize}
\end{esempio}

\subsubsection{Valutazione di un'espressione}

Sia $e$ un'espressione booleana sulle variabili $x_1, \dots, x_n$ e sia
$(v_1, \dots, v_n) \in \{0,1\}^n$ una configurazione delle variabili. La
\emph{valutazione} di $e$ rispetto a $(v_1, \dots, v_n)$ è il procedimento con
cui si sostituisce a ciascuna variabile il suo valore e, applicando le regole
dell'algebra di commutazione, si ottiene un valore $v \in \{0,1\}$ per
l'espressione. Ripetendo il procedimento per tutte le configurazioni si ottiene
una tabella di verità: \emph{ad ogni espressione corrisponde dunque una
funzione}.

\begin{esempio}
Sia $e = z(y + x) + xy$. La valutazione per $(1,0,0)$ è
\[
  e(100) = 0 \cdot (0 + 1) + 1 \cdot 0 = 0 .
\]
La valutazione completa determina la tabella di verità della funzione booleana
su $x, y, z$:
\begin{center}
\begin{tabular}{ccc|c}
$x$ & $y$ & $z$ & $e$ \\ \midrule
0 & 0 & 0 & 0 \\
0 & 0 & 1 & 0 \\
0 & 1 & 0 & 0 \\
0 & 1 & 1 & 1 \\
1 & 0 & 0 & 0 \\
1 & 0 & 1 & 1 \\
1 & 1 & 0 & 1 \\
1 & 1 & 1 & 1
\end{tabular}
\end{center}
\end{esempio}

\subsubsection{Espressioni e funzioni}

Ad una singola espressione corrisponde una singola funzione; ad una funzione
possono invece corrispondere molte espressioni.

\begin{definizione}[Espressioni equivalenti]
Due espressioni $e_1, e_2$ definite sullo stesso insieme di variabili si dicono
\emph{equivalenti}, e si scrive $e_1 \equiv e_2$, se corrispondono alla stessa
funzione.
\end{definizione}

Le trasformazioni eseguite su un'espressione di partenza $e_0$ usando i
postulati e le proprietà dell'algebra booleana portano a espressioni
equivalenti.

\subsubsection{Espressioni a confronto}

Nei contesti in cui si usa un'algebra booleana è fondamentale rappresentare le
funzioni con \emph{buone} espressioni. Un possibile criterio di bontà consiste
nel minimizzare il numero di letterali: si individua una prima espressione che
rappresenti la funzione e si cerca poi un'espressione equivalente migliore (con
meno letterali). Espressioni di partenza tipiche da manipolare per minimizzare i
letterali sono le espressioni in \emph{forma canonica}. Le forme canoniche di
base sono due: la prima espressa come somma di prodotti (SOP), la seconda come
prodotto di somme (POS); entrambe si ottengono dal teorema di espansione di
Shannon.

\subsection{Forme canoniche}

\subsubsection{Teorema di espansione di Shannon}

\begin{teorema}[Espansione di Shannon]
Sia $f \colon B^n \to B$ una funzione booleana di $n$ variabili. Allora per
ogni $(x_1, \dots, x_n) \in B^n$
\[
  f(x_1, \dots, x_n) = x_1' \cdot f(0, x_2, \dots, x_n)
                     + x_1 \cdot f(1, x_2, \dots, x_n)
\]
e dualmente
\[
  f(x_1, \dots, x_n) = \bigl(x_1' + f(1, x_2, \dots, x_n)\bigr)
                       \cdot \bigl(x_1 + f(0, x_2, \dots, x_n)\bigr).
\]
\end{teorema}

I termini $f(0, x_2, \dots, x_n)$ e $f(1, x_2, \dots, x_n)$ si chiamano
\emph{cofattori} e si indicano spesso con $f_{x_1'}$ e $f_{x_1}$.

\begin{esempio}[Applicazione iterata, forma somma]
Per una funzione di due variabili $f(x,y)$:
\begin{align*}
f(x,y) &= x' \cdot f(0,y) + x \cdot f(1,y)\\
       &= x' \bigl(y' \cdot f(0,0) + y \cdot f(0,1)\bigr)
          + x \bigl(y' \cdot f(1,0) + y \cdot f(1,1)\bigr)\\
       &= x'y' \cdot f(00) + x'y \cdot f(01) + xy' \cdot f(10) + xy \cdot f(11).
\end{align*}
In generale, l'applicazione iterata del teorema a una funzione di $n$ variabili
permette di riscriverla come somma di $2^n$ termini:
\begin{align*}
f(x_1,\dots,x_n) = {}& x_1' \cdots x_{n-1}' x_n' \cdot f(0\dots00)\\
 + {}& x_1' \cdots x_{n-1}' x_n \cdot f(0\dots01)\\
 + {}& \cdots\\
 + {}& x_1 \cdots x_{n-1} x_n \cdot f(1\dots11).
\end{align*}
\end{esempio}

\begin{esempio}[Applicazione iterata, forma prodotto]
Applicando la forma duale a $f(x,y)$:
\begin{align*}
f(x,y) &= \bigl(x' + f(1,y)\bigr)\bigl(x + f(0,y)\bigr)\\
       &= \bigl(x' + [(y' + f(11))(y + f(10))]\bigr)
          \bigl(x + [(y' + f(01))(y + f(00))]\bigr)\\
       &= (x' + y' + f(11))(x' + y + f(10))(x + y' + f(01))(x + y + f(00)).
\end{align*}
In generale si ottiene il prodotto di $2^n$ termini:
\begin{align*}
f(x_1,\dots,x_n) = {}& (x_1 + \dots + x_{n-1} + x_n + f(0\dots00))\\
 \cdot {}& (x_1 + \dots + x_{n-1} + x_n' + f(0\dots01))\\
 \cdot {}& \cdots\\
 \cdot {}& (x_1' + \dots + x_{n-1}' + x_n' + f(1\dots11)).
\end{align*}
\end{esempio}

\subsubsection{Prima forma canonica: mintermini}

\begin{definizione}[Mintermine]
Un \emph{mintermine} su $x_1, \dots, x_n$ è un termine prodotto
$y_1 \cdot y_2 \cdots y_n$ dove $y_i = x_i$ oppure $y_i = x_i'$, per ogni
$i = 1, \dots, n$.

Dato $\bar{v} = (v_1, \dots, v_n) \in \{0,1\}^n$ si costruisce il mintermine
$m_{\bar{v}} = y_1 \cdot y_2 \cdots y_n$ con $y_i = x_i$ se $v_i = 1$ e
$y_i = x_i'$ se $v_i = 0$ (il mintermine \emph{corrispondente} alla
configurazione $\bar{v}$).
\end{definizione}

Il mintermine $m_{\bar{v}}$ è un'espressione che vale $1$ in corrispondenza
dell'\emph{unica} configurazione di ingresso $\bar{v}$.

\begin{esempio}
Sia $\bar{v} = (101)$ una configurazione per le variabili $\{x,y,z\}$. Il
mintermine associato è $m_{\bar{v}} = xy'z$, che assume valore $1$ soltanto se
$x = 1$, $y = 0$ e $z = 1$, cioè per la sola configurazione $(101)$.
\end{esempio}

\begin{definizione}[Prima forma canonica (\emph{Sum of Products}, SOP)]
Data una funzione booleana $f(x_1, \dots, x_n)$, la sua rappresentazione in
\emph{prima forma canonica} è la somma di tutti i mintermini su
$x_1, \dots, x_n$ associati alle configurazioni per cui $f$ vale $1$.
\end{definizione}

L'applicazione iterata del teorema di Shannon porta alla prima forma canonica.
In pratica, per ricavarla basta guardare le righe della tabella di verità dove
$f$ vale $1$.

\begin{esempio}
Sia $f$ la funzione su $x, y, z$ che vale $1$ se e solo se esattamente due
variabili valgono $1$ (tabella dell'esempio all'inizio della lezione). I mintermini associati alle
configurazioni per cui $f = 1$ sono:
\begin{itemize}
    \item $x'yz$, dalla quarta riga ($011$);
    \item $xy'z$, dalla sesta riga ($101$);
    \item $xyz'$, dalla settima riga ($110$).
\end{itemize}
Quindi la prima forma canonica di $f$ è
\[
  f = x'yz + xy'z + xyz' .
\]
\end{esempio}

\subsubsection{Seconda forma canonica: maxtermini}

\begin{definizione}[Maxtermine]
Un \emph{maxtermine} su $x_1, \dots, x_n$ è un termine somma
$y_1 + y_2 + \dots + y_n$ dove $y_i = x_i$ oppure $y_i = x_i'$, per ogni
$i = 1, \dots, n$.

Data la configurazione $\bar{v} = (v_1, \dots, v_n) \in \{0,1\}^n$ si costruisce
il maxtermine $M_{\bar{v}} = y_1 + y_2 + \dots + y_n$ con $y_i = x_i$ se
$v_i = 0$ e $y_i = x_i'$ se $v_i = 1$ (il maxtermine \emph{corrispondente} alla
configurazione $\bar{v}$).
\end{definizione}

Il maxtermine $M_{\bar{v}}$ è un'espressione che vale $0$ in corrispondenza
dell'\emph{unica} configurazione di ingresso $\bar{v}$.

\begin{esempio}
Per $\bar{v} = (101)$ su $\{x,y,z\}$ il maxtermine associato è
$M_{\bar{v}} = x' + y + z'$, che assume valore $0$ soltanto se $x = 1$,
$y = 0$ e $z = 1$.
\end{esempio}

\begin{definizione}[Seconda forma canonica (\emph{Product of Sums}, POS)]
Data una funzione booleana $f(x_1, \dots, x_n)$, la sua rappresentazione in
\emph{seconda forma canonica} è il prodotto di tutti i maxtermini su
$x_1, \dots, x_n$ corrispondenti alle configurazioni per cui $f$ vale $0$.
\end{definizione}

L'applicazione iterata del teorema di Shannon nella sua forma duale porta alla
seconda forma canonica. In pratica, per ricavarla basta guardare le righe della
tabella di verità dove $f$ vale $0$.

\begin{esempio}
Per la stessa funzione (esattamente due variabili valgono $1$) i maxtermini
associati alle configurazioni per cui $f = 0$ sono:
\begin{itemize}
    \item $x + y + z$, dalla prima riga ($000$);
    \item $x + y + z'$, dalla seconda riga ($001$);
    \item $x + y' + z$, dalla terza riga ($010$);
    \item $x' + y + z$, dalla quinta riga ($100$);
    \item $x' + y' + z'$, dall'ultima riga ($111$).
\end{itemize}
Quindi la seconda forma canonica di $f$ è
\[
  f = (x + y + z)(x + y + z')(x + y' + z)(x' + y + z)(x' + y' + z') .
\]
\end{esempio}

\subsection*{Riferimenti}

Capitolo 2 del libro di testo: C.\ Bolchini et al., \emph{Reti Logiche},
Maggioli Editore. Lezione 2: sezione 2.3; lezione 3: sezioni 2.4 e 2.5.

\subsection*{Esercizi per casa}

\begin{esercizio}
Definire la tabella di verità della funzione booleana $f \colon B^3 \to B$ che
assume valore $1$ soltanto se almeno due bit in ingresso hanno valore $0$.
\end{esercizio}
\begin{esercizio}
Rappresentare $f$ mediante un'espressione booleana in prima forma canonica.
\end{esercizio}
\begin{esercizio}
Rappresentare $f'$ mediante un'espressione booleana in seconda forma canonica.
\end{esercizio}

\end{document}
